\documentclass[10pt,a4paper]{article}
\usepackage[T1]{fontenc}
\usepackage{lmodern,microtype}
\usepackage[margin=18mm]{geometry}
\usepackage{amsmath,amssymb,stmaryrd,mathpartir,booktabs,array}
\usepackage[dvipsnames]{xcolor}
\usepackage[colorlinks,linkcolor=MidnightBlue,urlcolor=MidnightBlue]{hyperref}
\hypersetup{pdftitle={PomPom open signatures: rules and semantics},
  pdfauthor={PomPom Language}}
\setlength{\parindent}{0pt}
\setlength{\parskip}{4pt}
\setlength{\emergencystretch}{3em}
\newcommand{\Set}{\mathsf{Set}}
\newcommand{\op}[1]{\mathsf{#1}}
\newcommand{\one}{\mathbf{1}}
\newcommand{\unit}{\ast}
\newcommand{\code}[1]{{}^{\prime}\!#1}
\newcommand{\dprod}{\mathbin{\code{\times}}}
\newcommand{\interp}[2]{\llbracket #1\rrbracket\,#2}
\newcommand{\Def}{\op{Def}}
\newcommand{\Fam}{\op{Fam}}
\newcommand{\cl}[2]{\op{close}\,#1\,#2}
\newcommand{\roll}{\op{in}}
\newcommand{\jd}[2]{\Gamma\vdash #1:#2}
\newcommand{\subt}[3]{\Gamma\vdash #1\sqsubseteq #2\rightsquigarrow #3}
\newcommand{\dsub}[4]{\Gamma\vdash #1\preceq_{#3}#2\rightsquigarrow #4}
\newcommand{\dead}[3]{\Gamma\vdash #1\ \op{dead}_{#2}\rightsquigarrow #3}
\newcommand{\synth}[3]{\Gamma\vdash #1\Rightarrow #2\rightsquigarrow #3}
\newcommand{\checkt}[3]{\Gamma\vdash #1\Leftarrow #2\rightsquigarrow #3}
\newcommand{\rn}[1]{{\normalfont\ttfamily\footnotesize\detokenize{#1}}}
\newcommand{\ruleof}[3]{\inferrule*[right={\rn{#1}}]{#2}{#3}}
\newcommand{\rows}[1]{\langle #1\rangle}
\newcommand{\abort}{\op{abort}}
\newcommand{\unroll}{\op{unroll}}
\newcommand{\allh}{\op{iAll}}
\newcommand{\hyps}{\op{hyps}}
\newcommand{\disp}{\op{dispatch}}
\newcommand{\pos}[1]{\overline{#1}}
\newcommand{\rto}{\longmapsto}
\newcommand{\red}{\leadsto}
\newcommand{\obs}{\approx}
\newcommand{\imp}{\ \Longrightarrow\ }
\newcommand{\ex}{\par\smallskip\textit{Example.}\ }
\newcommand{\tg}[1]{\code{\mathsf{#1}}}
\newcommand{\spi}{\pi}
\newcommand{\nilE}{\op{nilE}}
\newcommand{\consE}[2]{\op{consE}\,#1\,#2}
\newcommand{\suc}[1]{1{+}\,#1}
\newcommand{\valid}{\op{valid}}
\newcommand{\pz}{\pi_0}
\newcommand{\po}{\pi_1}
\newcommand{\dg}{^{\dagger}}
\newcommand{\DM}{[DM]}
\newcommand{\ruleofd}[3]{\inferrule*[right={\rn{#1}$\dg$}]{#2}{#3}}
\newcommand{\Part}[2]{\par\bigskip\hrule\smallskip
 {\Large\bfseries Part #1.\quad #2}\par\smallskip\hrule\medskip}
\begin{document}
\begin{center}
{\Large\bfseries PomPom open signatures: typing rules and semantics}\\[3pt]
{\small A condensed reference for \texttt{progress/OpenSignatures*.v}
\quad\textbullet\quad 25 November 2026}
\end{center}

The core calculus is the type theory of Dagand and McBride,
\emph{Elaborating Inductive Definitions} (arXiv:1210.6390, Section~1; cited
as \DM), \emph{extended}. \textbf{Part~I} gives it: the rules that come from
\DM{} are written in that paper's notation and cited by figure, and every
rule that is ours, or that departs from the paper's, is marked $\dg$ and
explained where it stands. The extensions are explicit universe levels and
cumulativity, the empty code $\code{\bot}$, computation rules for the
induction hypotheses, and \emph{parameterized closure}.
\textbf{Part~II} is elaboration: a source language with named constructors,
exhaustive cases and subtyping, translated into core terms. It adds no core
rule; each of its judgments outputs a core term, and its soundness theorems
say that output is typed by Part~I.

Labels in \rn{typewriter} are the Rocq constructor or theorem names
(\texttt{progress/OpenSignatures*.v}). All stated theorems are proved in
Rocq; the one refuted statement is marked. The rules, definitions, theorem
statements and examples below are restated in Rocq and proved from the names
they cite in \texttt{progress/OpenSignaturesDocCheck\{Core,Elab,Rel\}.v}
(\texttt{make -C progress check-docs}); names beginning with \rn{ex_} are
example lemmas of those files.

\textbf{Conventions.} Notation is that of \DM: $(x{:}S)\to T$ and
$(x{:}S)\times T$ are the dependent function and pair types, written $S\to T$
and $S\times T$ when $x$ does not occur in $T$; $\pz$, $\pi_1$ are the
projections; $\one$ is the unit type and $\unit$ its element. $k,j$ are
universe levels, $s$ is a string, and $\pos{n}$ is the position
$\suc{(\cdots(\suc{0}))}$ with $n$ successors. Binders are stable numeric
IDs, contexts are finite maps, and substitution $t[s/x]$ is capture
avoiding. An ambient index type $I:\Set_0$ is fixed and suppressed: in Rocq
it is an explicit argument of every operator on descriptions and of the
$\op{dead}$, $\preceq$ and handler judgments of Part~II.


\Part{I}{The core calculus}
%----------------------------------------------------------------------
\section{Syntax}
\[
\begin{array}{@{}r@{\ }c@{\ }l@{\quad}l@{}}
t,S,T &::=& x\mid\Set_k\mid(x{:}S)\to T\mid\lambda x.\,t\mid t\,s
  \mid(x{:}S)\times T\mid(s,t)\mid\pz\,t\mid\po\,t
  \mid\one\mid\unit &\text{type theory}\\
&\mid&\op{UId}\mid\code{s}\mid\op{EnumU}\mid\nilE\mid\consE{t}{E}\mid\op{EnumT}\,E
  \mid 0\mid\suc{n}\mid\spi^k\,E\,P\mid\op{switch}^k\,E\,P\,b\,x
  &\text{enumerations}\\
&\mid&\op{IDesc}\,I\mid\code{\op{var}}\,i\mid\code{\one}\mid\code{\bot}
  \mid A\dprod B\mid\code{\Pi}\,S\,T\mid\code{\Sigma}\,S\,T
  \mid\code{\sigma}\,E\,T\mid\interp{D}{X} &\text{descriptions}\\
&\mid&\allh\,D\,X\,t\,P\mid\hyps\,D\,X\,P\,h\,t
  \mid\mu\,D\mid\roll\,t\mid\op{ind}\,D\,P\,s\,i\,t &\text{fix-point, induction}\\
&\mid&\cl{F}{G}\mid\op{closeCase}^k\,F\,G\,i\,Q\,b\,t
  \mid\op{closeInd}\,G\,P\,s\,F\,i\,t &\text{closure}
\end{array}
\]
Abbreviations are introduced below, together with the operators they
abbreviate.

%----------------------------------------------------------------------
\section{Typing and computation rules}\label{sec:core}
The judgments are $\Gamma\vdash\valid$ and $\Gamma\vdash t:T$ (\rn{wf},
\rn{typing}). The groups below follow Figures 1, 2 and 4 of \DM, each with
its purpose, its typing rules, and its \emph{root computation} rules
$t\rto u$ (\rn{root_step}); Section~\ref{sec:dyn} closes $\rto$ into
evaluation, full reduction and conversion $\equiv$.

\textbf{Reading the cited rules.} An unmarked rule is the paper's rule. In
this calculus it is read with three standing differences, which are not
repeated rule by rule:
\begin{itemize}\setlength{\itemsep}{0pt}\setlength{\topsep}{2pt}
\item \emph{Levels.} Where \DM{} writes an unsubscripted $\Set$, the rule
 here is at $\Set_0$.
\item \emph{Presuppositions.} What a conclusion mentions must be well
 formed, and here that is a premise of the rule: the function or pair type in
 the rules for application, pairing and projections; $t:\op{UId}$ and
 $E:\op{EnumU}$ in the rules for $0$ and $\suc{n}$; $I:\Set_0$ in every rule
 that mentions $\op{IDesc}\,I$; and $D:I\to\op{IDesc}\,I$, $i:I$ in the rule
 for $\roll$.
\item \emph{Annotations.} $\lambda x.\,t$ and $(s,t)$ carry no type
 annotation (\DM: $\lambda x{:}S.\,t$ and $(s,t)_{x.T}$), and contexts hold no
 definitions.
\end{itemize}
\textbf{Signatures.} As in \DM, an operator with several arguments is given
by a signature $\op{op}:(x_1{:}S_1)\cdots(x_n{:}S_n)\to T$. It abbreviates
the typing rule with one premise $\Gamma\vdash s_i:S_i$ per argument,
earlier arguments being substituted for $x_1,x_2,\ldots$, and conclusion
$\Gamma\vdash\op{op}\,s_1\cdots s_n:T$. Operators are term formers: they are
always fully applied.

\subsection{Type theory \DM{} Fig.~1}\label{sec:tt}
The base is a Martin-L\"of type theory with $\Sigma$-types and $\Pi$-types
and a hierarchy of universes $\Set_k$.

\textbf{(a) Context validity.}
\begin{mathpar}
\ruleof{wf_nil}{}{\vdash\valid}
\and
\ruleof{wf_cons}{\Gamma\vdash\valid\\\Gamma\vdash S:\Set_k\\x\notin\Gamma}
 {\Gamma;x{:}S\vdash\valid}
\end{mathpar}
\textbf{(b) Judgmental equality.} \DM{} specifies equality by three typed
equations. Here the same three are root computation rules,
\[
(\lambda x.\,t)\,s\rto t[s/x]\qquad\pz\,(s,t)\rto s\qquad\po\,(s,t)\rto t
\]
and equality is the \emph{untyped} conversion $\equiv$ that all root
computation rules generate (Section~\ref{sec:dyn}). It also contains
$\eta$ for functions$\dg$, which is not a root computation.

\textbf{(c) Typing judgments.} In any valid $\Gamma$: \
$\Set_k:\Set_{k+1}$, \ $\one:\Set_k$, \ $\unit:\one$ \
(\rn{ty_sort}, \rn{ty_unitT}, \rn{ty_unit}).
\begin{mathpar}
\ruleof{ty_var}{\Gamma;x{:}S;\Delta\vdash\valid}{\Gamma;x{:}S;\Delta\vdash x:S}
\and
\ruleofd{ty_conv}{\Gamma\vdash s:S\\\Gamma\vdash T:\Set_k\\S\equiv T}{\Gamma\vdash s:T}
\and
\ruleofd{ty_alpha}{\Gamma\vdash t:T\\t\equiv_\alpha t'}{\Gamma\vdash t':T}
\\
\ruleofd{ty_sigma}{\Gamma\vdash S:\Set_j\\\Gamma;x{:}S\vdash T:\Set_k}
 {\Gamma\vdash(x{:}S)\times T:\Set_{\max(j,k)}}
\and
\ruleof{ty_pair}{\Gamma\vdash s:S\\\Gamma\vdash t:T[s/x]}
 {\Gamma\vdash(s,t):(x{:}S)\times T}
\\
\ruleof{ty_fst}{\Gamma\vdash p:(x{:}S)\times T}{\Gamma\vdash\pz\,p:S}
\and
\ruleof{ty_snd}{\Gamma\vdash p:(x{:}S)\times T}{\Gamma\vdash\po\,p:T[\pz\,p/x]}
\\
\ruleofd{ty_pi}{\Gamma\vdash S:\Set_j\\\Gamma;x{:}S\vdash T:\Set_k}
 {\Gamma\vdash(x{:}S)\to T:\Set_{\max(j,k)}}
\and
\ruleofd{ty_lam}{\Gamma\vdash(x{:}S)\to T:\Set_k\\\Gamma;x{:}S\vdash t:T}
 {\Gamma\vdash\lambda x.\,t:(x{:}S)\to T}
\\
\ruleof{ty_app}{\Gamma\vdash f:(x{:}S)\to T\\\Gamma\vdash s:S}
 {\Gamma\vdash f\,s:T[s/x]}
\end{mathpar}
In every rule that extends the context, $x\notin\Gamma$. Marked rules:
\rn{ty_conv} has an untyped $\equiv$ with $T$ well formed, where \DM{} has
the typed premise $\Gamma\vdash S\equiv T:\Set_k$; the two formation rules
take the maximum of two levels, where \DM{} uses one level $k$ throughout;
\rn{ty_lam}, whose $\lambda$ has no annotation, takes the well-formedness of
the function type where \DM{} has $\Gamma\vdash S:\Set_k$; \rn{ty_alpha} is
needed because binders are named IDs.

\textbf{(d) Cumulativity$\dg$.} \DM{} writes $\Set_k$ ``implicitly assuming
cumulativity of universes''. Here it is explicit: a type in $\Set_j$ is also
in every larger universe, and a function may be used at a function type
with a smaller domain and a larger codomain universe.
\begin{mathpar}
\ruleofd{ty_cumul}{\Gamma\vdash t:\Set_j\\j\le k}{\Gamma\vdash t:\Set_k}
\and
\ruleofd{ty_cumul_fun}{\Gamma\vdash f:(x{:}S)\to T\\\Gamma\vdash(x{:}S)\to T:\Set_j\\
 \Gamma\vdash(x{:}S')\to T':\Set_k\\S'\le_u S\\T\le_u T'}
 {\Gamma\vdash f:(x{:}S')\to T'}
\\
\ruleofd{ul_refl}{}{S\le_u S}
\and
\ruleofd{ul_sort}{j\le k}{\Set_j\le_u\Set_k}
\and
\ruleofd{ul_pi}{S'\le_u S\\T\le_u T'}{(x{:}S)\to T\le_u(x{:}S')\to T'}
\end{mathpar}
Cumulativity is not part of $\equiv$.
\ex If $f:\one\to\Set_0$ then $f:\one\to\Set_1$, and if $g:\Set_1\to\Set_0$
then $g:\Set_0\to\Set_0$, both by \rn{ty_cumul_fun}; yet
$\Set_0\not\equiv\Set_1$ (\texttt{OpenSignaturesEtaRegression.v}). Before
this rule was added, $\lambda x.\,f\,x$ had the type $\one\to\Set_1$ and its
$\eta$-reduct $f$ did not
(\texttt{proof-archive/2026-09-27-before-function-cumulativity}).

\subsection{Universe of enumerations \DM{} Fig.~2}\label{sec:enum}
This universe defines finite collections of labels. Labels are introduced
through the type $\op{UId}$ and are used to define finite sets through the
universe $\op{EnumU}$. To index a specific element of such a set one writes
an $\op{EnumT}$ code. To eliminate a finite set, the small $\Pi$-type $\spi$
builds a \emph{lookup tuple} that maps every label of the enumeration to a
value, and the eliminator $\op{switch}$ performs the lookup.

\textbf{(a) Tags.} A tag $\code{s}$ is written with a string $s$; nothing
computes on it.
\begin{mathpar}
\ruleof{ty_uid}{\Gamma\vdash\valid}{\Gamma\vdash\op{UId}:\Set}
\and
\ruleof{ty_tag}{\Gamma\vdash\valid}{\Gamma\vdash\code{s}:\op{UId}}
\end{mathpar}
\textbf{(b) Enumeration.}
\begin{mathpar}
\ruleof{ty_enumu}{\Gamma\vdash\valid}{\Gamma\vdash\op{EnumU}:\Set}
\and
\ruleof{ty_nile}{\Gamma\vdash\valid}{\Gamma\vdash\nilE:\op{EnumU}}
\and
\ruleof{ty_conse}{\Gamma\vdash t:\op{UId}\\\Gamma\vdash E:\op{EnumU}}
 {\Gamma\vdash\consE{t}{E}:\op{EnumU}}
\end{mathpar}
\textbf{(c) Index.}
\begin{mathpar}
\ruleof{ty_enumt}{\Gamma\vdash E:\op{EnumU}}{\Gamma\vdash\op{EnumT}\,E:\Set}
\and
\ruleof{ty_zero}{\Gamma\vdash\valid}{\Gamma\vdash 0:\op{EnumT}\,(\consE{t}{E})}
\and
\ruleof{ty_succ}{\Gamma\vdash n:\op{EnumT}\,E}
 {\Gamma\vdash\suc{n}:\op{EnumT}\,(\consE{t}{E})}
\end{mathpar}
\textbf{(d) Lookup tuples and elimination.} \DM{} gives $\spi$ and
$\op{switch}$ by their types, with motives $P:\op{EnumT}\,E\to\Set$. Here
they carry the universe level $k$ of the motive$\dg$:
\[
\begin{array}{@{}r@{\ :\ }l@{\qquad}l@{}}
\spi^k & (E{:}\op{EnumU})\,(P{:}\op{EnumT}\,E\to\Set_k)\to\Set_k & \text{\rn{ty_epi}}\\
\op{switch}^k & (E{:}\op{EnumU})\,(P{:}\op{EnumT}\,E\to\Set_k)\to
  \spi^k\,E\,P\to(x{:}\op{EnumT}\,E)\to P\,x & \text{\rn{ty_switch}}
\end{array}
\]
and their computation rules are
\[
\begin{array}{@{}r@{\ \rto\ }l@{}}
\spi^k\,\nilE\,P & \one\\
\spi^k\,(\consE{t}{E})\,P & P\,0\times\spi^k\,E\,(\lambda n.\,P\,(\suc{n}))\\
\op{switch}^k\,(\consE{t}{E})\,P\,(b,b_s)\,0 & b\\
\op{switch}^k\,(\consE{t}{E})\,P\,(b,b_s)\,(\suc{n}) &
  \op{switch}^k\,E\,(\lambda n.\,P\,(\suc{n}))\,b_s\,n
\end{array}
\]
So $\spi^k\,E\,P$ has one component of type $P\,\pos{n}$ per position.
Neither $\op{UId}$ nor $\op{EnumU}$ has an eliminator of its own: $\spi$ is
the only recursion over a code, and $\op{switch}$ the only use of a position.

\textbf{Notation} (\DM{} Examples 1 and 2). A set of labels is written by
enumerating them, and a map out of it by giving a value for each:
\[
\begin{array}{@{}r@{\ :=\ }l@{}}
\{\tg{a},\tg{b},\tg{c}\} & \consE{\tg{a}}{(\consE{\tg{b}}{(\consE{\tg{c}}{\nilE})})}\ :\ \op{EnumU}\\
\{\tg{a}\mapsto e_a,\ \tg{b}\mapsto e_b,\ \tg{c}\mapsto e_c\} &
 \lambda x.\,\op{switch}^k\,\{\tg{a},\tg{b},\tg{c}\}\,P\,(e_a,(e_b,(e_c,\unit)))\,x
 \ :\ (x{:}\op{EnumT}\,\{\tg{a},\tg{b},\tg{c}\})\to P\,x
\end{array}
\]
with $P$ and $k$ left implicit (\rn{ex_lookup_map}). \DM{} applies
$\op{switch}$ to three arguments; here it is fully applied under a
$\lambda$.
\ex For $E=\{\tg{nil},\tg{cons}\}$ the positions are $0$ and $\suc{0}$, and
$\spi^k\,E\,P\equiv P\,0\times(P\,(\suc{0})\times\one)$: a lookup tuple is
$(b_0,(b_1,\unit))$. Looking up the second label,
$\{\tg{nil}\mapsto b_0,\tg{cons}\mapsto b_1\}\,(\suc{0})$, takes three root
steps: $\beta$, then
$\op{switch}^k\,E\,P\,(b_0,(b_1,\unit))\,(\suc{0})\rto
\op{switch}^k\,\{\tg{cons}\}\,(\lambda n.\,P\,(\suc{n}))\,(b_1,\unit)\,0\rto b_1$.

\textbf{Empty type$\dg$.} The empty set of labels gives
$\bot:=\op{EnumT}\,\nilE$, with no positions, and
$\abort^k\,A\,z:=\op{switch}^k\,\nilE\,(\lambda\_.A)\,\unit\,z$ eliminates it
into any $A:\Set_k$ (\rn{abort_typing}, Section~\ref{sec:meta}).

\subsection{Universe of descriptions \DM{} Fig.~4(a)}\label{sec:desc}
To define a datatype one gives its \emph{code}, describing its signature
functor: which data a node stores and where its recursive children sit.
Only the indexed universe of \DM{} is used here (their Fig.~4, not the
non-indexed $\op{Desc}$ of Fig.~3). Children are indexed by $I$: a code is a
$D:\op{IDesc}\,I$, and $\code{\op{var}}\,i$ marks a child at index $i$. The
interpretation $\interp{D}{X}$ turns the code into the type of a node's
payload, given a family $X$ for the children. Write $\Fam:=I\to\Set_0$ for
families and $\Def:=I\to\op{IDesc}\,I$ for definitions, which give a code at
each index (\rn{Family}, \rn{Def}).

\textbf{(a) Codes.}
\begin{mathpar}
\ruleofd{ty_idesc}{\Gamma\vdash I:\Set_0}{\Gamma\vdash\op{IDesc}\,I:\Set_1}
\and
\ruleof{ty_ivar}{\Gamma\vdash i:I}{\Gamma\vdash\code{\op{var}}\,i:\op{IDesc}\,I}
\and
\ruleof{ty_i1}{\Gamma\vdash\valid}{\Gamma\vdash\code{\one}:\op{IDesc}\,I}
\and
\ruleofd{ty_ibot}{\Gamma\vdash\valid}{\Gamma\vdash\code{\bot}:\op{IDesc}\,I}
\\
\ruleof{ty_iprod}{\Gamma\vdash A:\op{IDesc}\,I\\\Gamma\vdash B:\op{IDesc}\,I}
 {\Gamma\vdash A\dprod B:\op{IDesc}\,I}
\and
\ruleof{ty_ipi, ty_isig}{\Gamma\vdash S:\Set\\\Gamma\vdash T:S\to\op{IDesc}\,I}
 {\Gamma\vdash\code{\Pi}\,S\,T,\ \ \code{\Sigma}\,S\,T:\op{IDesc}\,I}
\\
\ruleof{ty_ichoice}{\Gamma\vdash E:\op{EnumU}\\\Gamma\vdash T:\op{EnumT}\,E\to\op{IDesc}\,I}
 {\Gamma\vdash\code{\sigma}\,E\,T:\op{IDesc}\,I}
\end{mathpar}
Marked rules: \DM{} has $\op{IDesc}\,I:\Set$ for $I:\Set$, where the type of
codes is here one universe up, as their $\op{Desc}:\Set_1$; and
$\code{\bot}$, the code of the empty payload, is added.

\textbf{Interpretation.} A code is interpreted by the type former of
the same name; $\code{\op{var}}\,i$ is the child at $i$.
\[
\interp{(D{:}\op{IDesc}\,I)}{(X{:}I\to\Set)}:\Set\qquad\text{\rn{ty_interp}}
\]
\[
\begin{array}{@{}r@{\ \rto\ }l@{\qquad}r@{\ \rto\ }l@{}}
\interp{\code{\op{var}}\,i}{X} & X\,i &
\interp{A\dprod B}{X} & \interp{A}{X}\times\interp{B}{X}\\
\interp{\code{\one}}{X} & \one &
\interp{\code{\Pi}\,S\,T}{X} & (s{:}S)\to\interp{T\,s}{X}\\
\interp{\code{\bot}}{X} & \bot\quad\dg &
\interp{\code{\Sigma}\,S\,T}{X} & (s{:}S)\times\interp{T\,s}{X}\\
\multicolumn{2}{c}{} &
\interp{\code{\sigma}\,E\,T}{X} & (e{:}\op{EnumT}\,E)\times\interp{T\,e}{X}
\end{array}
\]
A datatype definition starts with a finite choice of constructors: a
$\code{\sigma}$ over the set of their names, with a map giving each one's
payload. A constructor whose payload is $\code{\bot}$ is \emph{hidden}: its
label is still in the set, but $\interp{\code{\bot}}{X}\rto\bot$ and no closed
term has that type, so it can never be built (\rn{ex_no_bot_payload}).
\ex Natural numbers (\DM{} Example~3, here with index type $I=\one$):
\[
\op{NatD}:=\lambda\_.\,\code{\sigma}\,\{\tg{zero},\tg{suc}\}\,
 \{\tg{zero}\mapsto\code{\one},\ \tg{suc}\mapsto\code{\op{var}}\,\unit\}\ :\ \Def.
\]
A payload in $\interp{\op{NatD}\,\unit}{X}$ is a position paired with that
constructor's data: $(0,\unit)$, or $(\suc{0},n)$ with $n:X\,\unit$
(\rn{ex_nat_def}, \rn{ex_nat_payloads}).
Interpretation computes under binders too:
$\interp{\code{\Sigma}\,A\,(\lambda\_.\code{\op{var}}\,\unit)}{X}\equiv A\times X\,\unit$
is one root step, then $\beta$ under the binder of the pair type (which
evaluation does not do, but $\equiv$ does; Section~\ref{sec:dyn}).

\subsection{Induction hypotheses$\dg$}\label{sec:hyps}
Induction needs, for a payload $xs$, one hypothesis for each recursive child
inside $xs$. In \DM{} the function $\op{iAll}$ ``computes the inductive
hypothesis'', and its definition is left to Chapman et al.\ (2010). Here
$\allh$ and $\hyps$ are operators with their own rules: $\allh$ computes the
type of the collection of hypotheses, and $\hyps$ builds the collection from
a function $h$ that gives the result at every child. For a family $X$,
$\op{Mot}(X):=((i{:}I)\times X\,i)\to\Set_0$ is the type of motives and
$\op{Rec}(X,P):=(i{:}I)\to(y{:}X\,i)\to P\,(i,y)$ the type of $h$
(\rn{motive}, \rn{recursive_method}).
\[
\begin{array}{@{}r@{\ :\ }l@{\qquad}l@{}}
\allh & (D{:}\op{IDesc}\,I)\,(X{:}\Fam)\,(xs{:}\interp{D}{X})\,(P{:}\op{Mot}(X))\to\Set_0
 & \text{\rn{ty_iall}}\\
\hyps & (D{:}\op{IDesc}\,I)\,(X{:}\Fam)\,(P{:}\op{Mot}(X))\,(h{:}\op{Rec}(X,P))\,
 (xs{:}\interp{D}{X})\to\allh\,D\,X\,xs\,P & \text{\rn{ty_hyps}}
\end{array}
\]
\textbf{Computation.} Both follow the code $D$ and the shape of the payload.
Write $\mathcal A_D(xs)=\allh\,D\,X\,xs\,P$ and
$\mathcal H_D(xs)=\hyps\,D\,X\,P\,h\,xs$.
\begin{center}
\renewcommand{\arraystretch}{1.25}
\begin{tabular}{@{}l@{\qquad}l@{\qquad}l@{\qquad}l@{}}
\toprule
$D$ & $xs$ & $\mathcal A_D(xs)\rto$ & $\mathcal H_D(xs)\rto$\\
\midrule
$\code{\op{var}}\,i$ & $y$ & $P\,(i,y)$ & $h\,i\,y$\\
$\code{\one}$ & $\unit$ & $\one$ & $\unit$\\
$\code{\bot}$ & $z$ & $\one$ & $\unit$\\
$A\dprod B$ & $(a,b)$ &
 $\mathcal A_{A}(a)\times\mathcal A_{B}(b)$ &
 $(\mathcal H_{A}(a),\mathcal H_{B}(b))$\\
$\code{\Pi}\,S\,T$ & $f$ &
 $(s{:}S)\to\mathcal A_{T\,s}(f\,s)$ & $\lambda s.\,\mathcal H_{T\,s}(f\,s)$\\
$\code{\Sigma}\,S\,T$ & $(s,y)$ &
 $\mathcal A_{T\,s}(y)$ & $\mathcal H_{T\,s}(y)$\\
$\code{\sigma}\,E\,T$ & $(e,y)$ &
 $\mathcal A_{T\,e}(y)$ & $\mathcal H_{T\,e}(y)$\\
\bottomrule
\end{tabular}
\end{center}
\ex For the two payloads of $\op{NatD}$: the $\tg{zero}$ payload $(0,\unit)$
has no child, so its hypotheses have type $\one$ and are $\unit$; the
$\tg{suc}$ payload $(\suc{0},n)$ has the one child $n$, so
$\allh\,(\op{NatD}\,\unit)\,X\,(\suc{0},n)\,P$ evaluates to $P\,(\unit,n)$ and
$\hyps\,(\op{NatD}\,\unit)\,X\,P\,h\,(\suc{0},n)$ to $h\,\unit\,n$
(\rn{ex_nat_hyps}).

\subsection{Fix-point \DM{} Fig.~4(b),(c)}\label{sec:mu}
The codes are interpreted to build the least fix-point: $\mu D$ is the
family whose values at index $i$ are $\roll\,xs$ for a payload $xs$ of $D\,i$
whose children are again in $\mu D$. Its argument is a function from indices
to codes. Its eliminator is induction. (\DM{} write $\mu^i$ and
$\op{iinduction}$.)

\textbf{(b) Fix-point.}
\begin{mathpar}
\ruleof{ty_mui}{\Gamma\vdash D:I\to\op{IDesc}\,I}{\Gamma\vdash\mu D:I\to\Set}
\and
\ruleof{ty_in_mui}{\Gamma\vdash xs:\interp{D\,i}{(\mu D)}}{\Gamma\vdash\roll\,xs:\mu D\,i}
\end{mathpar}
\textbf{(c) Elimination principle.}
\[
\begin{array}{@{}r@{\ :\ }l@{}}
\op{ind} & (D{:}I\to\op{IDesc}\,I)\,(P{:}((i{:}I)\times\mu D\,i)\to\Set)\to{}\\
 & \bigl((i{:}I)\,(xs{:}\interp{D\,i}{(\mu D)})\to
   \allh\,(D\,i)\,(\mu D)\,xs\,P\to P\,(i,\roll\,xs)\bigr)\to
   (i{:}I)\,(x{:}\mu D\,i)\to P\,(i,x)
 \qquad\text{\rn{ty_ind}}
\end{array}
\]
Its computation rule calls the step with the payload and its hypotheses,
built by recursive calls:
\[
\op{ind}\,D\,P\,s\,i\,(\roll\,xs)\rto
 s\,i\,xs\,\bigl(\hyps\,(D\,i)\,(\mu D)\,P\,
   (\lambda j\,y.\,\op{ind}\,D\,P\,s\,j\,y)\,xs\bigr)
\]
\ex $\op{Nat}:=\mu\,\op{NatD}\,\unit$, with $\op{zero}:=\roll\,(0,\unit)$ and
$\op{suc}\,n:=\roll\,(\suc{0},n)$ (\rn{ex_nat_values}). Induction on
$\op{suc}\,n$ calls the step with the payload $(\suc{0},n)$ and the
hypothesis $h\,\unit\,n$ of the previous example, where $h$ is the recursive
call: the hypothesis is $\op{ind}\,\op{NatD}\,P\,s\,\unit\,n$.

\subsection{Parameterized closure$\dg$}\label{sec:close}
This subsection is the extension; nothing in it comes from \DM. In $\mu D$
one definition describes every layer of a value. $\cl{F}{G}$ separates two
roles: $F$ describes the \emph{outer} layer of a value, and $G$ describes
every layer below it. So $\cl{G}{G}$ is the full datatype of $G$, and
$\cl{F}{G}$ is a type whose root node has the shape $F$ and whose children
are full values of $G$. Write
\[
C_{F,G}(i):=\cl{F}{G}\,i,\qquad X_G:=\cl{G}{G}
\]
(\rn{CloseAt}, \rn{carrier}): a value of $C_{F,G}(i)$ is $\roll\,xs$ with
payload $xs:\interp{F\,i}{X_G}$.

\textbf{(a) Formation and constructor.} $F$ and $G$ are unrelated: there is
no inclusion premise.
\begin{mathpar}
\ruleofd{ty_close}{\Gamma\vdash F:\Def\\\Gamma\vdash G:\Def}{\Gamma\vdash\cl{F}{G}:\Fam}
\and
\ruleofd{ty_in_close}{\Gamma\vdash F:\Def\\\Gamma\vdash G:\Def\\\Gamma\vdash i:I\\
 \Gamma\vdash xs:\interp{F\,i}{X_G}}
 {\Gamma\vdash\roll\,xs:C_{F,G}(i)}
\end{mathpar}
\textbf{(b) Case analysis.} A value is opened one layer; the motive may live
in any universe.
\[
\begin{array}{@{}r@{\ :\ }l@{}}
\op{closeCase}^k & (F,G{:}\Def)\,(i{:}I)\,(Q{:}C_{F,G}(i)\to\Set_k)\to{}\\
 & \bigl((xs{:}\interp{F\,i}{X_G})\to Q\,(\roll\,xs)\bigr)\to
 (x{:}C_{F,G}(i))\to Q\,x\qquad\text{\rn{ty_close_case}}
\end{array}
\]
\[
\op{closeCase}^k\,F\,G\,i\,Q\,b\,(\roll\,xs)\rto b\,xs
\]
In particular
$\unroll_{F,G,i}\,x:=\op{closeCase}^0\,F\,G\,i\,(\lambda\_.\interp{F\,i}{X_G})\,(\lambda y.y)\,x$
returns the payload:
$\unroll_{F,G,i}\,(\roll\,xs)\rto(\lambda y.y)\,xs\rto xs$
(\rn{unroll_typing}, \rn{unroll_beta}, Section~\ref{sec:meta}).

\textbf{(c) Elimination principle.} The scrutinee has outer definition $F$,
but its children are values of $X_G$, with outer definition $G$. The motive
therefore quantifies over the outer definition, and
$P^G:=\lambda z.\,P\,G\,(\pz\,z)\,(\po\,z)$ is its instance at the children
(\rn{diagonal_motive}). Inductive motives live in $\Set_0$.
\[
\begin{array}{@{}r@{\ :\ }l@{}}
\op{closeInd} & (G{:}\Def)\,
 \bigl(P{:}(H{:}\Def)\to(j{:}I)\to C_{H,G}(j)\to\Set_0\bigr)\to{}\\
 & \bigl((H{:}\Def)\,(j{:}I)\,(xs{:}\interp{H\,j}{X_G})\to
   \allh\,(H\,j)\,X_G\,xs\,P^G\to P\,H\,j\,(\roll\,xs)\bigr)\to{}\\
 & (F{:}\Def)\,(i{:}I)\,(x{:}C_{F,G}(i))\to P\,F\,i\,x
 \qquad\text{\rn{ty_close_ind}}
\end{array}
\]
\[
\op{closeInd}\,G\,P\,s\,F\,i\,(\roll\,xs)\rto
 s\,F\,i\,xs\,\bigl(\hyps\,(F\,i)\,X_G\,P^G\,
   (\lambda j\,y.\,\op{closeInd}\,G\,P\,s\,G\,j\,y)\,xs\bigr)
\]
The recursive call switches the outer definition from $F$ to $G$. There is
\emph{no} rule unfolding a $\op{close}$ type, and $C_{G,G}$ is not identified
with $\mu\,G$; the constructor $\roll$ is shared by both.
\ex Positive numbers keep both labels of $\op{NatD}$ and hide $\tg{zero}$:
\[
\op{PosD}:=\lambda\_.\,\code{\sigma}\,\{\tg{zero},\tg{suc}\}\,
 \{\tg{zero}\mapsto\code{\bot},\ \tg{suc}\mapsto\code{\op{var}}\,\unit\}.
\]
$C_{\op{NatD},\op{NatD}}(\unit)$ is the natural numbers again, and
$C_{\op{PosD},\op{NatD}}(\unit)$ is the positive ones. Its values are
$\roll\,(\suc{0},n)$ with $n$ an arbitrary natural number of
$C_{\op{NatD},\op{NatD}}(\unit)$: the \emph{same terms} as the successors
there, at the same position (\rn{ex_pos_values}). No closed value uses
$\tg{zero}$, whose payload would have type $\bot$: every closed term of the
type evaluates to such a successor (\rn{ex_pos_only_suc}). Hiding
constructors of the outer definition while keeping the recursive one is how
Part~II obtains subtypes.

%----------------------------------------------------------------------
\section{Evaluation, reduction and conversion}\label{sec:dyn}
The root computations $\rto$ of Section~\ref{sec:core} generate three relations:
\begin{center}
\renewcommand{\arraystretch}{1.25}
\begin{tabular}{@{}l@{\quad}l@{\quad}p{112mm}@{}}
\toprule
$t\to u$ & \rn{step} & \emph{Evaluation.} $\rto$ under weak contexts
$\mathcal E::=[\,]\mid\mathcal E\,a\mid(\mathcal E,b)\mid(a,\mathcal E)\mid
\pz\,\mathcal E\mid\po\,\mathcal E\mid\suc{\mathcal E}\mid\roll\,\mathcal E$,
the code $E$ of $\spi$, the $E$, tuple and scrutinee of $\op{switch}$,
the description of $\interp{-}{X}$, the description and data argument of
$\allh$/$\hyps$, and the scrutinee of
$\op{ind}$, $\op{closeCase}$, $\op{closeInd}$. Never under $\lambda$ or in
an argument. $\to^*$ is \rn{eval}; \rn{next_step}/\rn{run} is a
deterministic strategy for it.\\
$t\red u$ & \rn{reduction} & \emph{Full reduction.} $\rto$ in every
position, including under binders, plus
$\lambda x.\,f\,x\red f$ $(x\notin\op{FV}(f))$.\\
$t\equiv u$ & \rn{conv} & \emph{Conversion.} Least equivalence that is a
congruence for every term former and contains $\to$, $\eta$ and
$\equiv_\alpha$.\\
\bottomrule
\end{tabular}
\end{center}
\textbf{Values} (\rn{value}): $\Set_k$, function and pair types, $\lambda$, pairs,
$\one$, $\unit$, $\op{UId}$, $\code{s}$, $\op{EnumU}$, $\nilE$, $\consE{t}{E}$,
$\op{EnumT}\,E$, $0$, $\suc{n}$, $\op{IDesc}\,I$, all description codes,
$\mu D$, $\mu D\,i$, $\cl{F}{G}$, $\cl{F}{G}\,i$, and $\roll\,x$.
\ex $t=\lambda x.\,(\lambda y.y)\,x$ is a value, so it has no $\to$ step,
but $t\red\lambda x.x$ ($\beta$ under the binder) and $t\red\lambda y.y$
($\eta$).

%----------------------------------------------------------------------
\section{Observational equivalence}\label{sec:obs}
Let $\mathbb B:=\op{EnumT}\,\{\tg{true},\tg{false}\}$.
A closing substitution $\gamma$ for $\Gamma$ (\rn{closing}) assigns to each
$x{:}A$ a closed $u$ with $\varnothing\vdash u:A\gamma$.
\[
\begin{array}{@{}r@{\ }c@{\ }l@{}}
t\obs^{\varnothing}_A u &:\Leftrightarrow&
 \varnothing\vdash t,u:A\ \wedge\ \forall\,(x{:}A\vdash c:\mathbb B),\ v\in\{0,\suc{0}\}.\
 \ c[t/x]\to^* v\iff c[u/x]\to^* v\\[2pt]
\Gamma\vdash t\obs_A u &:\Leftrightarrow&
 \Gamma\vdash t,u:A\ \wedge\ \forall\gamma\text{ closing }\Gamma.\ \
 t\gamma\obs^{\varnothing}_{A\gamma}u\gamma
\end{array}
\]
(\rn{closed_observational_eq}, \rn{observational_eq}.) $\obs$ is an
equivalence, contains typed conversion, is (on closed terms) a congruence
for application and typed contexts, is invariant under $\equiv$ on the type, and is equivalently
tested by closed functions $A\to\mathbb B$
(\rn{observational_conversion}, \rn{closed_observational_congruence},
\rn{closed_observational_context}, \rn{closed_observational_application},
\rn{observational_type_conversion}, \rn{closed_observational_function_tests}).
\ex For a variable $x:C_{F,G}(i)$, the terms $\roll\,(\unroll_{F,G,i}\,x)$
and $x$ have no common reduct, so they are not convertible; but every
closing substitution replaces $x$ by a term evaluating to some $\roll\,y$,
after which $\unroll$ computes, so they are $\obs$
(\rn{close_roll_unroll}).

%----------------------------------------------------------------------
\section{Core metatheory}\label{sec:meta}
Unless noted, statements are in the \texttt{Theorems} module; the two about
full reduction are in \texttt{ComputationPreservation} and
\texttt{EtaPolymorphism}. The \texttt{Audit} module prints all their
assumptions.

\newcommand{\thm}[2]{\rn{#1} & $#2$\\}
\renewcommand{\arraystretch}{1.3}
\textbf{Structure.}\par
\begin{tabular}{@{}l@{\quad}l@{}}
\thm{context_validity}{\jd{t}{A}\imp\Gamma\vdash\valid}
\thm{type_correctness}{\jd{t}{A}\imp\exists k.\ \jd{A}{\Set_k}}
\thm{weakening}{\jd{t}{A},\ \jd{B}{\Set_k},\ x\notin\Gamma\imp\Gamma;x{:}B\vdash t:A}
\thm{substitution}{x\notin\Gamma,\ \Gamma;x{:}A\vdash t:B,\ \jd{u}{A}\imp\jd{t[u/x]}{B[u/x]}}
\thm{closing_substitution}{\jd{t}{A},\ \gamma\text{ closing }\Gamma\imp\varnothing\vdash t\gamma:A\gamma}
\thm{conversion_substitution}{t\equiv u\imp t[s/x]\equiv u[s/x]}
\end{tabular}

\textbf{Reduction.}\par
\begin{tabular}{@{}l@{\quad}l@{}}
\thm{preservation}{\jd{t}{A},\ t\to u\imp\jd{u}{A}\quad(\text{and for }\to^*:\ \text{\rn{preservation_eval}})}
\thm{named_computation_preservation}{\jd{t}{A},\ t\red_{\beta}u\imp\jd{u}{A}
 \quad(\red_\beta:\ \rto\text{ in every position, no }\eta)}
\thm{progress}{\varnothing\vdash t:A\imp\op{value}(t)\ \vee\ \exists u.\ t\to u}
\thm{normalization}{\jd{t}{A}\imp t\text{ is strongly normalizing for }\red}
\thm{normalization_eval}{\varnothing\vdash t:A\imp\exists v.\ t\to^*v\wedge\op{value}(v)}
\thm{confluence}{\jd{t}{A},\ t\red^*u,\ t\red^*v\imp\exists w\equiv_\alpha w'.\ u\red^*w\wedge v\red^*w'}
\thm{conversion_joinability}{\jd{t}{A},\ \jd{u}{A},\ t\equiv u\imp\exists w\equiv_\alpha w'.\ t\red^*w\wedge u\red^*w'}
\thm{consistency}{\neg\exists t.\ \varnothing\vdash t:\bot}
\thm{canonical_forms_close}{\varnothing\vdash v:C_{F,G}(i),\ \op{value}(v)\imp
 \exists x.\ v=\roll\,x\wedge\varnothing\vdash x:\interp{F\,i}{X_G}}
\end{tabular}

Both also hold for raw terms, with the typing premises dropped
(\rn{raw_confluence}, \rn{raw_conversion_joinability}).

\textbf{Refuted} (\rn{full_preservation_refuted}): preservation for $\red$
with $\eta$ fails:
\[
T=(A{:}\Set_0)\to A\to A,\quad f=(\lambda z.\lambda B.z)(\lambda x.x):\qquad
\vdash\lambda A.\,f\,A:T,\qquad\lambda A.\,f\,A\red f,\qquad\nvdash f:T.
\]
Unannotated $\lambda$ lets the type of $z$ depend on $A$ only inside the
$\eta$-expansion. $\eta$ therefore lives in $\equiv$ and $\red$ but not in
$\to$; safety uses $\to$, and normalization/confluence of $\red$ need no
preservation for it.

\textbf{Derived core rules.} Write $\op{CI}$ for the premise
$\jd{F,G}{\Def}\wedge\jd{i}{I}$ (\rn{close_input}).\par
\begin{tabular}{@{}l@{\quad}l@{}}
\thm{abort_typing}{\jd{A}{\Set_k},\ \jd{z}{\bot}\imp\jd{\abort^k\,A\,z}{A}}
\thm{unroll_typing}{\op{CI},\ \jd{x}{C_{F,G}(i)}\imp\jd{\unroll_{F,G,i}\,x}{\interp{F\,i}{X_G}}}
\thm{unroll_beta}{\unroll_{F,G,i}(\roll\,x)\to^*x}
\thm{close_roll_unroll}{\op{CI},\ \jd{x}{C_{F,G}(i)}\imp\Gamma\vdash\roll\,(\unroll_{F,G,i}\,x)\obs_{C_{F,G}(i)}x
 \quad(\text{not a conversion})}
\end{tabular}

\section{Core lemmas and proof structure}\label{sec:lemmas}
\renewcommand{\arraystretch}{1.25}
\begin{tabular}{@{}>{\raggedright}p{26mm}@{\quad}p{140mm}@{}}
\toprule
Binding &
Substitution respects $\equiv_\alpha$ and composes up to $\equiv_\alpha$;
$\op{FV}(t\sigma)=\bigcup_{v\in\op{FV}(t)}\op{FV}(\sigma v)$; a binder may be
renamed to any ID not free in the body; typed terms are scoped by their
context (\rn{substitute_alpha}, \rn{substitute_compose},
\rn{free_vars_substitute}, \rn{alpha_rename_lam/pi/sigma},
\rn{typing_scoped}, \rn{typing_fresh_not_free}). Typing depends only on the
context as a map, hence exchange (\rn{typing_ctx_equal},
\rn{ctx_equal_exchange}).\\
Nameless transport &
The named calculus is encoded into a de Bruijn calculus
(\texttt{nameless/DB*.v}); typing is preserved and reflected
(\rn{typing_encoding}, \rn{typing_reflection}), and reduction commutes with
encoding. Weakening, substitution, type correctness, preservation,
confluence and normalization are proved there and transported back
(\rn{named_weakening}, \rn{named_substitution}, \rn{named_preservation},
\rn{raw_confluence}, \rn{named_full_normalization}).\\
Conversion inversion &
From joinability: $\equiv$ preserves rigid heads (type formers, codes, tags,
positions, $\roll$; not $\lambda$ or pairs, because of $\eta$); equal tags
have equal strings; enumeration codes, $\code{\sigma}$ codes and function-type
domains are injective up to $\equiv$
(\rn{raw_head}, \rn{conv_tags}, \rn{conv_conse}, \rn{conv_choice},
\rn{pi_domain_conversion}); $\Set_0\not\equiv\Set_1$ and
$\Set_1\not\le_u\Set_0$.\\
Canonical forms &
A value at a close type is $\roll\,x$; at $\op{EnumT}\,E$ it is $0$ or
$\suc{n}$; at $\interp{\code{\sigma}\,E\,T}{X}$ it is a pair; a closed term
of $\op{EnumT}\,\{\code{s_0},\ldots,\code{s_{m-1}}\}$ evaluates to some
$\pos{n}$ with $n<m$; no value has
type $\bot$ (\rn{close_value_shape}, \rn{enum_value_shape},
\rn{choice_value_shape}, \rn{enum_normal_form}, \rn{bottom_no_value}).
Progress follows from these plus joinability
(\rn{progress_from_conversion}); consistency from normalization, preservation
and \rn{bottom_no_value}.\\
Normalization &
A cumulative reducibility-candidate model of the nameless calculus:
a cumulative tower of universes, descriptions as positive (monotone)
functors with positive fixed points for $\mu$ and $\op{close}$,
computability of every primitive operator, then the fundamental theorem
(\rn{semantic_fundamental}, \rn{typing_full_normalization}).\\
\bottomrule
\end{tabular}


\newpage
\Part{II}{Elaboration into the core}
%----------------------------------------------------------------------
\section{Elaboration: rows, dead branches, coercions}\label{sec:elab}
The source language has named constructors, exhaustive case and subtyping;
all three elaborate to core terms of Part~I. Every judgment below
\emph{outputs} a core term after $\rightsquigarrow$; none of them is used by
a core rule. These are relations, not algorithms.

\textbf{The idea: hide constructors with $\code{\bot}$.} A subtype of a
datatype keeps \emph{all} of its constructor labels and gives the
constructors it excludes the empty payload $\code{\bot}$. Three things then
follow from $\bot$ being empty. The programmer cannot build a hidden
constructor, because its payload would be a term of type $\bot$. A case
analysis need not handle it, because the missing branch is discharged by
$\abort$. And a coercion to a larger type never meets it, so the value passes
through unchanged, at the same position. The rules below implement this with
the judgment $\op{dead}$, which certifies that a payload type is empty, and
they relate two rows only when both list the same constructor names.

\subsection{Rows and signatures}
A datatype definition starts with a finite choice of constructors
(Section~\ref{sec:desc}). A \emph{row}
$r=[(\ell_0,D_0),\ldots,(\ell_{n-1},D_{n-1})]$ lists constructor names with
their payload descriptions; it denotes the $\code{\sigma}$ code over the set
of those names, with the map sending each name to its payload (the map is
at level $1$, with motive $\lambda\_.\op{IDesc}\,I$).
\[
\begin{array}{@{}l@{\qquad}l@{}}
E_r:=\{\code{\ell_0},\ldots,\code{\ell_{n-1}}\} &
p_r:=\{\code{\ell_0}\mapsto D_0,\ldots,\code{\ell_{n-1}}\mapsto D_{n-1}\}\\[2pt]
\rows{r}:=\code{\sigma}\,E_r\,p_r\quad(\text{\rn{row_code}}) &
\op{sig}(i.\,r):=\lambda i.\,\rows{r}\quad(\text{\rn{signature}})\\[2pt]
\multicolumn{2}{@{}l@{}}{\op{RowOK}_\Gamma(r)\ :\Leftrightarrow\
 \text{names of $r$ distinct}\ \wedge\ \forall(\ell,D)\in r.\ \jd{D}{\op{IDesc}\,I}
 \quad(\text{\rn{row_input}})}\\[2pt]
\multicolumn{2}{@{}l@{}}{\Gamma\vdash D\Downarrow r\ :\Leftrightarrow\
 \op{RowOK}_\Gamma(r)\ \wedge\ \jd{D}{\op{IDesc}\,I}\ \wedge\ D\equiv\rows{r}
 \quad(\text{\rn{row_view}})}
\end{array}
\]
\textbf{Names and positions.} The constructor names of a row become the
tags of $E_r$; core terms only ever compute with positions.
$\op{At}(E,s,n)$ (\rn{label_at}) says that name $s$ sits at position $n$ of
$E$: \ $\op{At}(\consE{\code{s}}{E},s,0)$ (\rn{at_here}); \
$\op{At}(E,s,n)\Rightarrow\op{At}(\consE{t}{E},s,\suc{n})$ (\rn{at_there}); closed
under $\equiv$ in $E$ and $n$ (\rn{at_convert}). Every entry satisfies
$\op{At}(E_r,\ell_n,\pos{n})$, and with distinct names the position of a name
is unique up to $\equiv$ (Section~\ref{sec:emeta}).

\textbf{Handlers}, one per row entry, are assembled by a dispatcher on the
tag:
\[
\begin{array}{@{}l@{}}
\disp^{k}_{r,X,Y}(h_s):=\lambda z.\,\bigl(\op{switch}^k\,E_r\,
  (\lambda c.\,\interp{p_r\,c}{X}\to Y)\,(h_0,(\ldots,\unit))\,(\pz\,z)\bigr)(\po\,z)
  \quad(\text{\rn{row_map}})\\[2pt]
\op{retag}_m:=\lambda y.\,(\pos{m},y)\qquad
\op{exfalso}^{k}_{Y}(d):=\lambda y.\,\abort^k\,Y\,(d\,y)
\end{array}
\]
\textbf{Running example.} Take $I=\one$, an element type $A:\Set_0$,
$K_A:=\code{\Sigma}\,A\,(\lambda\_.\code{\op{var}}\,\unit)$, and the rows and
signatures
\[
r_L:=[(\texttt{nil},\code{\one}),(\texttt{cons},K_A)],\quad
r_N:=[(\texttt{nil},\code{\bot}),(\texttt{cons},K_A)],\quad
L_A:=\op{sig}(\_.\,r_L),\quad N_A:=\op{sig}(\_.\,r_N),
\]
that is,
$L_A=\lambda\_.\,\code{\sigma}\,\{\tg{nil},\tg{cons}\}\,\{\tg{nil}\mapsto\code{\one},\ \tg{cons}\mapsto K_A\}$
and
$N_A=\lambda\_.\,\code{\sigma}\,\{\tg{nil},\tg{cons}\}\,\{\tg{nil}\mapsto\code{\bot},\ \tg{cons}\mapsto K_A\}$:
the same labels, with \texttt{nil} hidden. Then
\[
\op{List}\,A:=C_{L_A,L_A}(\unit),\qquad\op{NonEmpty}\,A:=C_{N_A,L_A}(\unit).
\]
Since $\interp{K_A}{X_{L_A}}\equiv A\times\op{List}\,A$, the values of
$\op{List}\,A$ are $\op{nil}=\roll\,(0,\unit)$ and $\roll\,(\pos{1},(a,x))$,
and the values of $\op{NonEmpty}\,A$ are $\roll\,(\pos{1},(a,x))$ only, with
the tail $x$ a full $\op{List}\,A$ in both cases. A non-empty list is the
same term as the corresponding list, and \texttt{cons} sits at position $1$
in both rows: $\op{At}(E_{r_L},\texttt{cons},\pos{1})$ and
$\op{At}(E_{r_N},\texttt{cons},\pos{1})$. Every closed term of
$\op{NonEmpty}\,A$ evaluates to a \texttt{cons} (\rn{ex_nonempty_only_cons}).
In the sources these are \rn{list_rows}, \rn{nonempty_rows}, \rn{list_type}
and \rn{nonempty_type}.

\subsection{Dead descriptions and description coercions}
Standing premises of every $\op{dead}$ and $\preceq$ rule:
$\jd{D,D'}{\op{IDesc}\,I}$ and $\jd{X}{\Fam}$ (the handler rules \rn{rh_*}
have no typing premise of their own).

\textbf{(a) Dead descriptions.} $\dead{D}{X}{d}$ certifies emptiness with a witness $d:\interp{D}{X}\to\bot$.
\begin{mathpar}
\ruleof{dead_bot}{D\equiv\code{\bot}}{\dead{D}{X}{\lambda z.z}}
\and
\ruleof{dead_nil}{D\equiv\code{\sigma}\,\nilE\,T}{\dead{D}{X}{\lambda z.\,\pz\,z}}
\and
\ruleof{dead_provided}{\jd{d}{\interp{D}{X}\to\bot}}{\dead{D}{X}{d}}
\\
\ruleof{dead_left, dead_right}{D\equiv D_0\dprod D_1\\\dead{D_m}{X}{d}}
 {\dead{D}{X}{\lambda z.\,d(\pi_m\,z)}}
\and
\ruleof{dead_choice}{\Gamma\vdash D\Downarrow r\\
 \forall j.\ \dead{D_j}{X}{d_j}}
 {\dead{D}{X}{\disp^{0}_{r,X,\bot}(d_s)}}
\end{mathpar}
In \rn{dead_left} and \rn{dead_right}, $m$ is $0$ or $1$.

\textbf{(b) Description coercions and handlers.} $\dsub{D}{D'}{X}{q}$
outputs $q:\interp{D}{X}\to\interp{D'}{X}$. With $Y:=\interp{D'}{X}$:
\begin{mathpar}
\ruleof{ds_conv}{D\equiv D'}{\dsub{D}{D'}{X}{\lambda z.z}}
\and
\ruleof{ds_dead}{\dead{D}{X}{d}}{\dsub{D}{D'}{X}{\op{exfalso}^{0}_{Y}(d)}}
\and
\ruleof{ds_rows}{\Gamma\vdash D\Downarrow r\\\Gamma\vdash D'\Downarrow r'\\
 \text{names}(r)=\text{names}(r')\\
 \forall j.\ r_j\hookrightarrow_{r'}h_j}
 {\dsub{D}{D'}{X}{\disp^{0}_{r,X,Y}(h_s)}}
\\
\ruleof{rh_live}{r'[m]=(\ell,D'')\\\interp{D_j}{X}\equiv\interp{D''}{X}}
 {(\ell,D_j)\hookrightarrow_{r'}\op{retag}_m}
\and
\ruleof{rh_dead}{\dead{D_j}{X}{d}}
 {(\ell,D_j)\hookrightarrow_{r'}\op{exfalso}^{0}_{Y}(d)}
\end{mathpar}
\rn{ds_rows} relates two rows only if they list the same constructor names.
Each constructor of the source is then either \emph{hidden}, with a dead
payload such as $\code{\bot}$, or \emph{live}, with the same payload type as
in the target up to conversion. Since the names agree and are distinct, the
$m$ of \rn{rh_live} is the position of the source entry itself: a live
constructor keeps its position (\rn{ex_live_keeps_position}).
\ex $\rows{r_N}\preceq_X\rows{r_L}\rightsquigarrow
\disp^0_{r_N,X,Y}([\op{exfalso}^0_Y(\lambda z.z),\op{retag}_1])$:
\texttt{nil} is dead by \rn{dead_bot}, and \texttt{cons} is live and stays at
position $1$ (\rn{ex_embed_nonempty_list}). Going from $r_L$ to $r_N$ would
need \texttt{nil}, with payload $\code{\one}$, to be dead or to have the
payload $\code{\bot}$; it is neither, since $\code{\one}$ is not dead in a
closable context (see the case example below). A row that drops the label,
$[(\texttt{cons},K_A)]$, is related to neither of them: its names differ.

\subsection{Type coercions}
$\subt{A}{B}{c}$ outputs $c:A\to B$ (\rn{sub}).
\begin{mathpar}
\ruleof{su_conv}{\jd{A}{\Set_j}\\\jd{B}{\Set_k}\\A\equiv B}{\subt{A}{B}{\lambda x.x}}
\and
\ruleof{su_trans}{\subt{A}{B}{c}\\\subt{B}{C}{d}}{\subt{A}{C}{\lambda x.\,d(c\,x)}}
\and
\ruleof{su_bottom}{\jd{A}{\Set_k}}{\subt{\bot}{A}{\lambda z.\,\abort^k\,A\,z}}
\\
\ruleof{su_pi}{y\notin\Gamma\\
 \jd{(x{:}A)\to B}{\Set_j}\\\jd{(y{:}A')\to B'}{\Set_k}\\
 \subt{A'}{A}{c}\\
 \Gamma;y{:}A'\vdash B[c\,y/x]\sqsubseteq B'\rightsquigarrow d}
 {\subt{(x{:}A)\to B}{(y{:}A')\to B'}{\lambda f.\,\lambda y.\,d\,(f\,(c\,y))}}
\\
\ruleof{su_close}{\jd{F,H,G}{\Def}\\\jd{i}{I}\\\dsub{F\,i}{H\,i}{X_G}{q}}
 {\subt{C_{F,G}(i)}{C_{H,G}(i)}{\lambda x.\,\roll\,(q\,(\unroll_{F,G,i}\,x))}}
\end{mathpar}
\rn{su_close} changes only the outer layer; the recursive definition $G$ is
the same on both sides, so a coercion is a single retagging, not a traversal.
\ex By \rn{su_close} and the previous example,
$\subt{\op{NonEmpty}\,A}{\op{List}\,A}{\op{toList}}$ with
$\op{toList}=\lambda x.\,\roll\,(\disp^0([\op{exfalso}^0(\lambda z.z),\op{retag}_1])\,(\unroll\,x))$
(\rn{nonempty_widening}), and
$\op{toList}\,(\roll\,(\pos{1},(a,x)))\to^*\roll\,(\pos{1},(a,x))$: the
coercion returns the value it was given. There is no coercion the other way
for a variable $A$: by \rn{subtyping_sound} it would be a function
$\op{List}\,A\to\op{NonEmpty}\,A$ uniform in $A$, and none exists
(\rn{no_uniform_list_downcast}, \rn{ex_no_list_to_nonempty_coercion}). By
\rn{su_pi}, with \rn{su_conv} on a codomain $B$ that does not mention the
argument (\rn{ex_function_coercion}),
$\subt{\op{List}\,A\to B}{\op{NonEmpty}\,A\to B}
{\lambda f.\,\lambda y.\,(\lambda x.x)\,(f\,(\op{toList}\,y))}$.

\subsection{Source terms: synthesis and checking}
\[
e::=\op{core}(t)\mid x\mid(e:A)\mid\lambda x.\,e\mid e\,e\mid(e,e)\mid
 \op{sig}(i.\,r)\mid\op{close}\,e\,e\mid\ell(e)\mid
 \op{case}_Q\ e\ \op{of}\ \{\ell_j(y_j)\Rightarrow e_j\}_j
\]
The rules are organised bidirectionally, as the elaboration of \DM{}
(their Section~2.1) is; the rules themselves are specific to this calculus.

\textbf{(a) Type synthesis.}
\begin{mathpar}
\ruleof{es_var}{\Gamma\vdash\valid\\\Gamma(x)=A}{\synth{x}{A}{x}}
\and
\ruleof{es_ann}{\jd{A}{\Set_k}\\\checkt{e}{A}{t}}{\synth{(e:A)}{A}{t}}
\and
\ruleof{es_app}{\synth{f}{(x{:}A)\to B}{t}\\\checkt{a}{A}{u}}
 {\synth{f\,a}{B[u/x]}{t\,u}}
\\
\ruleof{es_signature}{i\notin\Gamma\\\op{RowOK}_{\Gamma;i{:}I}(r)}
 {\synth{\op{sig}(i.\,r)}{\Def}{\lambda i.\,\rows{r}}}
\and
\ruleof{es_close}{\checkt{F}{\Def}{f}\\\checkt{G}{\Def}{g}}
 {\synth{\op{close}\,F\,G}{\Fam}{\cl{f}{g}}}
\\
\ruleof{es_case}{\jd{F,G}{\Def}\\\jd{i}{I}\\\synth{e}{C_{F,G}(i)}{x}\\
 \jd{Q}{\Set_k}\\\Gamma\vdash F\,i\Downarrow r\\
 \text{clause names distinct and}\subseteq\text{names}(r)\\
 \forall j.\ r_j\rhd h_j}
 {\synth{\op{case}_Q\ e\ \op{of}\ b_s}{Q}
   {\op{closeCase}^k\,F\,G\,i\,(\lambda\_.Q)\,(\disp^{k}_{r,X_G,Q}(h_s))\,x}}
\end{mathpar}
\textbf{(b) Type checking.}
\begin{mathpar}
\ruleof{ec_core}{\jd{t}{A}}{\checkt{\op{core}(t)}{A}{t}}
\and
\ruleof{ec_lam}{x\notin\Gamma\\\Gamma;x{:}A\vdash e\Leftarrow B\rightsquigarrow t}
 {\checkt{\lambda x.\,e}{(x{:}A)\to B}{\lambda x.\,t}}
\and
\ruleof{ec_pair}{\checkt{e}{A}{t}\\\checkt{e'}{B[t/x]}{u}}
 {\checkt{(e,e')}{(x{:}A)\times B}{(t,u)}}
\\
\ruleof{ec_conversion}{\synth{e}{A}{t}\\A\equiv B}{\checkt{e}{B}{t}}
\and
\ruleof{ec_target_conversion}{\checkt{e}{A}{t}\\A\equiv B}{\checkt{e}{B}{t}}
\and
\ruleof{ec_subsumption}{\synth{e}{A}{t}\\\subt{A}{B}{c}}{\checkt{e}{B}{c\,t}}
\\
\ruleof{ec_constructor}{\jd{F,G}{\Def}\\\jd{i}{I}\\\Gamma\vdash F\,i\Downarrow r\\
 r[n]=(\ell,D)\\\checkt{e}{\interp{D}{X_G}}{u}}
 {\checkt{\ell(e)}{C_{F,G}(i)}{\roll\,(\pos{n},u)}}
\end{mathpar}
\textbf{(c) Case clauses.} One handler per entry $(\ell,D)$ of the row.
\begin{mathpar}
\ruleof{cases_live}{(\ell(y)\Rightarrow e')\in b_s\\y\notin\Gamma\\
 \Gamma;y{:}\interp{D}{X_G}\vdash e'\Leftarrow Q\rightsquigarrow t}
 {(\ell,D)\rhd\lambda y.\,t}
\and
\ruleof{cases_dead}{\ell\notin\text{names}(b_s)\\\dead{D}{X_G}{d}}
 {(\ell,D)\rhd\op{exfalso}^{k}_{Q}(d)}
\end{mathpar}
\rn{es_app}, \rn{ec_lam}, \rn{ec_pair} also require their function or pair
type well formed, and the two conversion rules require $B$ well formed.
The expected type fixes the row, hence the position $\pos{n}$ of a named
constructor. A case must cover every constructor of the row; an omitted one
needs a dead witness. $\op{core}(t)$ only checks, so raw $\roll$ terms never
have a guessed type.
\ex The source term $\texttt{cons}((\op{core}(a),\op{core}(\op{nil})))$, a
constructor applied to a pair, has no type of its own. Checked against
$\op{List}\,A$ or against $\op{NonEmpty}\,A$ it elaborates to the same term
$\roll\,(\pos{1},(a,\op{nil}))$ (\rn{ex_cons_at_list},
\rn{ex_cons_at_nonempty}). The source term $\texttt{nil}(e)$ checks against
$\op{List}\,A$ but not against $\op{NonEmpty}\,A$: there $e$ would have to
check against $\interp{\code{\bot}}{X}$, and no closed $e$ does
(\rn{ex_hidden_nil_unconstructible}). The function
$\op{head}:=\lambda x.\,\op{case}_A\ x\ \op{of}\ \{\texttt{cons}(p)\Rightarrow\op{core}(\pz\,p)\}$
checks against $\op{NonEmpty}\,A\to A$ with the \texttt{nil} clause omitted,
because the payload of \texttt{nil} is dead (\rn{cases_dead} with
\rn{dead_bot}). It elaborates to
$\lambda x.\,\op{closeCase}^0\,N_A\,L_A\,\unit\,(\lambda\_.A)\,
(\disp^0_{r_N,X_{L_A},A}([\op{exfalso}^0_A(\lambda z.z),\ \lambda p.\,\pz\,p]))\,x$
(\rn{ex_head_elaborates}), and
$\op{head}\,(\roll\,(\pos{1},(a,\op{nil})))\to^*a$ (\rn{singleton_head} at
$A=\one$). Against
$\op{List}\,A\to A$ the same case needs a handler for \texttt{nil}, i.e.\ a
dead witness $d:\interp{\code{\one}}{X}\to\bot$. None exists in a context that
has a closing substitution: $d\,\unit$ would close to a term of $\bot$
(\rn{ex_no_dead_unit}).

%----------------------------------------------------------------------
\section{Elaboration metatheory}\label{sec:emeta}
\renewcommand{\arraystretch}{1.3}
\textbf{Rows.}\par
\begin{tabular}{@{}l@{\quad}l@{}}
\thm{row_code_typing}{\op{RowOK}_\Gamma(r)\imp\jd{\rows{r}}{\op{IDesc}\,I}}
\thm{signature_typing}{i\notin\Gamma,\ \op{RowOK}_{\Gamma;i{:}I}(r)\imp\jd{\op{sig}(i.\,r)}{\Def}}
\thm{row_position_identity}{r[n]=(\ell,D)\imp\op{At}(E_r,\ell,\pos{n})}
\thm{label_resolution_unique}{\text{names distinct},\ \jd{a,b}{\op{EnumT}\,E_r},\ \op{At}(E_r,\ell,a),\ \op{At}(E_r,\ell,b)\imp a\equiv b}
\thm{canonical_forms_named}{\varnothing\vdash t:C_{F,G}(i),\ \varnothing\vdash F\,i\Downarrow r\imp
 \exists n,\ell,D,x.\ r[n]=(\ell,D)\wedge t\to^*\roll\,(\pos{n},x)}
 & $\hfill\wedge\ \varnothing\vdash x:\interp{D}{X_G}$\\
\end{tabular}

\textbf{Elaboration soundness and coherence.}\par
\begin{tabular}{@{}l@{\quad}l@{}}
\thm{dead_sound}{\dead{D}{X}{d}\imp\jd{d}{\interp{D}{X}\to\bot}}
\thm{dead_close_uninhabited}{\op{CI}\text{ in }\varnothing,\ \varnothing\vdash F\,i\ \op{dead}_{X_G}\rightsquigarrow d\imp
 \neg\exists x.\ \varnothing\vdash x:C_{F,G}(i)}
\thm{description_subtyping_sound}{\dsub{D}{D'}{X}{q}\imp\jd{q}{\interp{D}{X}\to\interp{D'}{X}}}
\thm{subtyping_sound}{\subt{A}{B}{c}\imp\jd{c}{A\to B}}
\thm{synthesis_sound, checking_sound}{\synth{e}{A}{t}\ \text{or}\ \checkt{e}{A}{t}\imp\jd{t}{A}}
\thm{elaborated_evaluation}{\varnothing\vdash e\Leftarrow A\rightsquigarrow t\imp
 \exists v.\ t\to^*v\wedge\op{value}(v)\wedge\varnothing\vdash v:A}
\thm{coercion_coherence}{\subt{A}{B}{c},\ \subt{A}{B}{d}\imp\Gamma\vdash c\obs_{A\to B}d}
\thm{checking_coherence}{\checkt{e}{A}{t},\ \checkt{e}{A}{u}\imp\Gamma\vdash t\obs_{A}u}
\end{tabular}

Coherence is stated with $\obs$, not $\equiv$: two coercions may differ as
terms yet agree on all closed inputs (compare the example of
Section~\ref{sec:obs}).

Handler lists are sound as well. \rn{row_handlers_sound}: if
$\op{RowOK}_\Gamma(r)$, $\op{RowOK}_\Gamma(r')$, $\jd{X}{\Fam}$,
$\jd{Y}{\Set_0}$, $Y\equiv\interp{\rows{r'}}{X}$ and
$r_j\hookrightarrow_{r'}h_j$ for all $j$, then each
$\jd{h_j}{\interp{D_j}{X}\to Y}$. \rn{case_handlers_sound}: if
$\op{RowOK}_\Gamma(r)$, $\jd{X}{\Fam}$, $\jd{Q}{\Set_k}$ and $r_j\rhd h_j$
for all $j$, then $\jd{(h_0,(\ldots,\unit))}
{\spi^k\,E_r\,(\lambda c.\,\interp{p_r\,c}{X}\to Q)}$.

\section{Elaboration lemmas and proof structure}
\renewcommand{\arraystretch}{1.25}
\begin{tabular}{@{}>{\raggedright}p{26mm}@{\quad}p{140mm}@{}}
\toprule
Rows &
$\disp(h_s)\,(\pos{n},x)\equiv h_n\,x$ and $\op{retag}_m\,x\equiv(\pos{m},x)$
(\rn{row_map_selected}, \rn{retag_handler_beta}); a closed payload of
$\rows{r}$ converts to some $(\pos{n},x)$ (\rn{row_payload_canonical}); two views of
one description agree on names per position with convertible payloads
(\rn{row_view_position_transport}); on a source entry with a closed
inhabitant every handler derivation picks some $\op{retag}_m$ with
$r'[m]$ of the same name (\rn{row_handlers_live_at}), and $m$ is unique when
the target names are distinct (\rn{row_position_unique}); under
\rn{ds_rows} it is the source position. Hence any two description coercions agree
pointwise up to $\equiv$ on closed inputs, and are the identity when
$D\equiv D'$ (\rn{description_coercions_closed_pointwise_coherence},
\rn{description_conversion_identity}).\\
Relational model &
Coherence needs more than pointwise agreement, so
\texttt{OpenSignaturesRel*.v} builds a binary logical relation on closed
terms: conversion-closed PERs, descriptions as relational functors with
least fixed points, a cumulative universe tower. \rn{rel_fundamental}: every
typed term is related to itself under related closing substitutions.
\rn{rel_adequacy}/\rn{rel_observational}: related terms are $\obs$.
Coercions are interpreted as semantic graphs (identity, empty source,
contravariant function, retag-by-name, roll). Two graphs between the same
types agree on related inputs, graphs are total, and the composite of two
graphs is contained in a graph (\rn{Gr_func}, \rn{Gr_total}, \rn{Gr_comp});
every
derivation of $\sqsubseteq$ or $\preceq$ realizes one (\rn{sub_real},
\rn{desc_real}). Functionality gives \rn{coercion_coherence}; a pairwise
induction on elaborations gives \rn{checking_coherence}.\\
\bottomrule
\end{tabular}

\section{The list example, as theorems}\label{sec:ex}
For the running example, with \texttt{nil} hidden by $\code{\bot}$
($\jd{A}{\Set_0}$ throughout):\par
\begin{tabular}{@{}l@{\quad}l@{}}
\thm{nonempty_widening}{\subt{\op{NonEmpty}\,A}{\op{List}\,A}{\op{toList}}}
\thm{nonempty_cons_typing}{\jd{a}{A},\ \jd{x}{\op{List}\,A}\imp
 \jd{\roll\,(\pos{1},(a,x))}{\op{NonEmpty}\,A}}
\thm{ex_nonempty_only_cons}{\varnothing\vdash t:\op{NonEmpty}\,A\imp
 \exists x.\ t\to^*\roll\,(\pos{1},x)}
\thm{nonempty_observers_typing}{\op{head}:\op{NonEmpty}\,A\to A,\quad
 \op{tail}:\op{NonEmpty}\,A\to\op{List}\,A}
\thm{no_uniform_list_downcast}{\neg\exists c.\ \varnothing\vdash c:(A{:}\Set_0)\to\op{List}\,A\to\op{NonEmpty}\,A}
\thm{self_restricted_list_empty}{\varnothing\vdash A:\Set_0\imp\neg\exists x.\ \varnothing\vdash x:C_{N_A,N_A}(\unit)}
\end{tabular}

$\op{head}$ and $\op{tail}$ are cases whose \texttt{nil} clause is
discharged by the dead payload. The last line is the list whose tail must
itself be non-empty: it has no closed value. The same $\op{cons}$ code is
reused in a tree signature (\rn{reused_tree_cons_typing}).
\end{document}
